% Options for packages loaded elsewhere
\PassOptionsToPackage{unicode}{hyperref}
\PassOptionsToPackage{hyphens}{url}
\PassOptionsToPackage{dvipsnames,svgnames,x11names}{xcolor}
%
\documentclass[
]{article}
\usepackage{amsmath,amssymb}
\usepackage{iftex}
\ifPDFTeX
  \usepackage[T1]{fontenc}
  \usepackage[utf8]{inputenc}
  \usepackage{textcomp} % provide euro and other symbols
\else % if luatex or xetex
  \usepackage{unicode-math} % this also loads fontspec
  \defaultfontfeatures{Scale=MatchLowercase}
  \defaultfontfeatures[\rmfamily]{Ligatures=TeX,Scale=1}
\fi
\usepackage{lmodern}
\ifPDFTeX\else
  % xetex/luatex font selection
\fi
% Use upquote if available, for straight quotes in verbatim environments
\IfFileExists{upquote.sty}{\usepackage{upquote}}{}
\IfFileExists{microtype.sty}{% use microtype if available
  \usepackage[]{microtype}
  \UseMicrotypeSet[protrusion]{basicmath} % disable protrusion for tt fonts
}{}
\makeatletter
\@ifundefined{KOMAClassName}{% if non-KOMA class
  \IfFileExists{parskip.sty}{%
    \usepackage{parskip}
  }{% else
    \setlength{\parindent}{0pt}
    \setlength{\parskip}{6pt plus 2pt minus 1pt}}
}{% if KOMA class
  \KOMAoptions{parskip=half}}
\makeatother
\usepackage{xcolor}
\usepackage[margin=1in]{geometry}
\usepackage{color}
\usepackage{fancyvrb}
\newcommand{\VerbBar}{|}
\newcommand{\VERB}{\Verb[commandchars=\\\{\}]}
\DefineVerbatimEnvironment{Highlighting}{Verbatim}{commandchars=\\\{\}}
% Add ',fontsize=\small' for more characters per line
\newenvironment{Shaded}{}{}
\newcommand{\AlertTok}[1]{\textcolor[rgb]{1.00,0.00,0.00}{\textbf{#1}}}
\newcommand{\AnnotationTok}[1]{\textcolor[rgb]{0.38,0.63,0.69}{\textbf{\textit{#1}}}}
\newcommand{\AttributeTok}[1]{\textcolor[rgb]{0.49,0.56,0.16}{#1}}
\newcommand{\BaseNTok}[1]{\textcolor[rgb]{0.25,0.63,0.44}{#1}}
\newcommand{\BuiltInTok}[1]{\textcolor[rgb]{0.00,0.50,0.00}{#1}}
\newcommand{\CharTok}[1]{\textcolor[rgb]{0.25,0.44,0.63}{#1}}
\newcommand{\CommentTok}[1]{\textcolor[rgb]{0.38,0.63,0.69}{\textit{#1}}}
\newcommand{\CommentVarTok}[1]{\textcolor[rgb]{0.38,0.63,0.69}{\textbf{\textit{#1}}}}
\newcommand{\ConstantTok}[1]{\textcolor[rgb]{0.53,0.00,0.00}{#1}}
\newcommand{\ControlFlowTok}[1]{\textcolor[rgb]{0.00,0.44,0.13}{\textbf{#1}}}
\newcommand{\DataTypeTok}[1]{\textcolor[rgb]{0.56,0.13,0.00}{#1}}
\newcommand{\DecValTok}[1]{\textcolor[rgb]{0.25,0.63,0.44}{#1}}
\newcommand{\DocumentationTok}[1]{\textcolor[rgb]{0.73,0.13,0.13}{\textit{#1}}}
\newcommand{\ErrorTok}[1]{\textcolor[rgb]{1.00,0.00,0.00}{\textbf{#1}}}
\newcommand{\ExtensionTok}[1]{#1}
\newcommand{\FloatTok}[1]{\textcolor[rgb]{0.25,0.63,0.44}{#1}}
\newcommand{\FunctionTok}[1]{\textcolor[rgb]{0.02,0.16,0.49}{#1}}
\newcommand{\ImportTok}[1]{\textcolor[rgb]{0.00,0.50,0.00}{\textbf{#1}}}
\newcommand{\InformationTok}[1]{\textcolor[rgb]{0.38,0.63,0.69}{\textbf{\textit{#1}}}}
\newcommand{\KeywordTok}[1]{\textcolor[rgb]{0.00,0.44,0.13}{\textbf{#1}}}
\newcommand{\NormalTok}[1]{#1}
\newcommand{\OperatorTok}[1]{\textcolor[rgb]{0.40,0.40,0.40}{#1}}
\newcommand{\OtherTok}[1]{\textcolor[rgb]{0.00,0.44,0.13}{#1}}
\newcommand{\PreprocessorTok}[1]{\textcolor[rgb]{0.74,0.48,0.00}{#1}}
\newcommand{\RegionMarkerTok}[1]{#1}
\newcommand{\SpecialCharTok}[1]{\textcolor[rgb]{0.25,0.44,0.63}{#1}}
\newcommand{\SpecialStringTok}[1]{\textcolor[rgb]{0.73,0.40,0.53}{#1}}
\newcommand{\StringTok}[1]{\textcolor[rgb]{0.25,0.44,0.63}{#1}}
\newcommand{\VariableTok}[1]{\textcolor[rgb]{0.10,0.09,0.49}{#1}}
\newcommand{\VerbatimStringTok}[1]{\textcolor[rgb]{0.25,0.44,0.63}{#1}}
\newcommand{\WarningTok}[1]{\textcolor[rgb]{0.38,0.63,0.69}{\textbf{\textit{#1}}}}
\setlength{\emergencystretch}{3em} % prevent overfull lines
\providecommand{\tightlist}{%
  \setlength{\itemsep}{0pt}\setlength{\parskip}{0pt}}
\setcounter{secnumdepth}{-\maxdimen} % remove section numbering
\setlength{\emergencystretch}{3em}
\ifLuaTeX
  \usepackage{selnolig}  % disable illegal ligatures
\fi
\usepackage{bookmark}
\IfFileExists{xurl.sty}{\usepackage{xurl}}{} % add URL line breaks if available
\urlstyle{same}
\hypersetup{
  colorlinks=true,
  linkcolor={black},
  filecolor={Maroon},
  citecolor={Blue},
  urlcolor={black},
  pdfcreator={LaTeX via pandoc}}

\author{}
\date{}

\begin{document}

\section{DDTA minimal BAE ontology candidate -
R1}\label{ddta-minimal-bae-ontology-candidate---r1}

\textbf{Status:} CANDIDATE / NOT ACCEPTED / BA1 OPEN

\textbf{Derived by:} BA1-T1

\textbf{Repository baseline reviewed:}
\texttt{83af68cb1a02a6b1e76f591d4c1519f9496be3b3}

\subsection{Candidate ontology}\label{candidate-ontology}

BA1-T1 proposes two first-class analytical identity families and no
domain-specific subtype split yet.

\subsubsection{\texorpdfstring{\texttt{BAReferent}}{BAReferent}}\label{bareferent}

An independently identifiable unit of methodology-neutral shared project
meaning that must be recognizable across one or more analytical
propositions, projections or governed baselines.

A \texttt{BAReferent} may carry a semantic kind/role such as structural
participant, external participant, capability, information/resource,
behavior/event, store, contract or other project meaning. BA1-T1 does
not promote those kinds to separate metaclasses.

\subsubsection{\texorpdfstring{\texttt{BAProposition}}{BAProposition}}\label{baproposition}

An independently identifiable methodology-neutral analytical assertion
about shared project meaning. It references one or more
\texttt{BAReferent} identities, may itself be targeted by other
propositions, and has enough identity to support origin/provenance,
review, diagnosis, reuse and change disposition.

Its content may express behavior, relation, dependency, responsibility,
correlation, state/condition, applicability, constraint, lifecycle or
failure semantics. BA2 will define the minimal relation/action
vocabulary; BA3 will define provenance and diagnostic materialization.

\subsection{Minimal picture}\label{minimal-picture}

\begin{Shaded}
\begin{Highlighting}[]
\NormalTok{BAReferent *}
\NormalTok{     \^{}}
\NormalTok{     | referenced by}
\NormalTok{     |}
\NormalTok{BAProposition *}

\NormalTok{BAProposition may also be the target of another BAProposition.}
\end{Highlighting}
\end{Shaded}

\texttt{BAE} is an umbrella term, not an additional required metaclass
in this candidate.

\subsection{Explicitly not accepted by
R1}\label{explicitly-not-accepted-by-r1}

No first-class split is accepted yet for:

\begin{itemize}
\tightlist
\item
  Participant / Actor;
\item
  Component / Capability;
\item
  Asset;
\item
  Information / Resource;
\item
  Behavior / Event / Transition;
\item
  Interface / Connection / Channel;
\item
  Store / Persistence;
\item
  Contract;
\item
  Boundary;
\item
  State / Mode / Context;
\item
  Dependency / Allocation;
\item
  Property / Constraint;
\item
  DataFlow / Interaction.
\end{itemize}

Document types (\texttt{MacroRequirement}, \texttt{Decision},
\texttt{FunctionalRequirement}, \texttt{SpecializedRequirement},
\texttt{SecurityRequirement}) remain source-layer governed concepts and
are not copied into Base Analysis as BAE types.

Diagnostic/unresolved, provenance/source locators and views are required
responsibilities but are deferred to BA3/BA4 material contracts rather
than proposed as domain BAE metaclasses.

\subsection{Falsification rule}\label{falsification-rule}

This candidate fails if a concrete governed corpus or method-neutral
consumer requires a deferred semantic kind to have independent
subtype-specific identity/invariants that cannot be represented honestly
with \texttt{BAReferent\ +\ BAProposition} without reconstructing
meaning outside the shared core.

\subsection{Next step}\label{next-step}

\texttt{BA1-T2\ -\ split-or-collapse\ pressure\ test\ of\ the\ two-family\ candidate}.

BA1 remains open. BA2 remains not started.

\end{document}
